% from ega4-14-fr.tex
\oldpage[IV-3]{210}

  \item Si $Y_0$ est l'unique composante irréductible de $Y$ contenant $y$, $\eta$
  son point générique, il existe une composante irréductible $Z$ de $X$ contenant
  $x$, dominant $Y_0$ et telle que
  $\dim_x(Z\cap X_y)=\dim(Z\cap X_\eta)$ (ce qui signifie que $Z$ est
  équidimensionnelle sur $Y_0$ au point $x$ (13.2.2)).
  \item[b$'$)] Pour tout voisinage ouvert $U$ de $x$ dans $X$ et toute
  générisation $y'$ de $y$, on a
  $\dim(U\cap X_{y'})\geq\dim_x(U\cap X_y)$.
\end{enumerate}
Si de plus $Y$ est localement noethérien, ces conditions sont aussi équivalentes à
la suivante~:
\begin{enumerate}[label=\emph{\alph*)}]
  \item[c)] $f$ est ouvert au point $x$.
\end{enumerate}
\end{theorem}

Notons d'abord que puisque $(\mathcal O_y)_{\mathrm{red}}$ est intègre, $y$
n'appartient qu'à une seule composante irréductible $Y_0$ de $Y$. Le fait que $b)$
et $b')$ sont équivalentes et que $a)$ implique $b')$ résulte de (14.3.13)~;
d'autre part, si $Y$ est localement noethérien, on a vu dans (14.2.3) que $c)$
implique $b)$. Il reste donc à montrer que lorsque $y$ est géométriquement
unibranche, $b)$ entraîne $a)$.

\begin{lemma}[14.4.1.1]
\label{IV.14.4.1.1-fr}
Soient $Y$ un préschéma, $Y'=Y[T_1,\ldots,T_n]$, $y'$ un point de $Y'$, $y$ son
image dans $Y$. Si $Y$ est géométriquement unibranche au point $y$, alors $Y'$ est
géométriquement unibranche au point $y'$.
\end{lemma}

En effet, comme pour tout $y\in Y$,
$\operatorname{Spec}(k(y)[T_1,\ldots,T_n])$ est géométriquement régulier sur $k(y)$
($\mathbf{0}_{\mathrm{IV}}$, 17.3.7), le morphisme structural $Y'\to Y$ est lisse
(6.8.1), et il suffit d'appliquer (11.3.14).

\begin{lemma}[14.4.1.2]
\label{IV.14.4.1.2-fr}
Soient $A$ un anneau local intègre unibranche, $B$ un anneau intègre contenant $A$
et entier sur $A$, $\mathfrak n$ un idéal premier de $B$ au-dessus de l'idéal
maximal $\mathfrak m$ de $A$. Alors le morphisme
$\operatorname{Spec}(B_{\mathfrak n})\to\operatorname{Spec}(A)$ est surjectif,
autrement dit, pour tout idéal premier $\mathfrak p$ de $A$, il existe un idéal
premier $\mathfrak q$ de $B$ tel que $\mathfrak q\subset\mathfrak n$ et
$\mathfrak q\cap A=\mathfrak p$.
\end{lemma}

Soient $K$ (resp. $L$) le corps des fractions de $A$ (resp. $B$), $A'$ la clôture
intégrale de $A$, $B'$ le sous-anneau de $L$ engendré par $A'$ et $B$, de sorte
qu'on a un diagramme commutatif d'injections canoniques
\[
  \xymatrix{
    B \ar[r] & B' \\
    A \ar[u] \ar[r] & A' \ar[u]
  }
\]
Comme $B'$ est entier sur $B$, il existe un idéal premier $\mathfrak n'$ de $B'$
tel que $\mathfrak n'\cap B=\mathfrak n$ (Bourbaki, \emph{Alg. comm.}, chap.~V,
\S~2, n\textsuperscript{o}~1, th.~1), et (pour la même raison)
$\operatorname{Spec}(A')\to\operatorname{Spec}(A)$ est surjectif. D'autre part,
comme $A$ est unibranche, $A'$ est un anneau \emph{local}, donc
$\mathfrak n'\cap A'$, qui est au-dessus de l'idéal maximal $\mathfrak m$ de $A$,
est nécessairement égal à l'unique idéal maximal $\mathfrak m'$ de $A'$. En vertu
du second théorème de Cohen--Seidenberg (\emph{loc. cit.}, \S~2,
n\textsuperscript{o}~4, th.~3) le morphisme
$\operatorname{Spec}(B'_{\mathfrak n'})\to\operatorname{Spec}(A')$ est surjectif,
donc il en est de même du composé
$\operatorname{Spec}(B'_{\mathfrak n'})\to\operatorname{Spec}(A')\to
\operatorname{Spec}(A)$~; mais ce morphisme est aussi le composé
$\operatorname{Spec}(B'_{\mathfrak n'})\to\operatorname{Spec}(B_{\mathfrak n})
\to\operatorname{Spec}(A)$, donc le morphisme
$\operatorname{Spec}(B_{\mathfrak n})\to\operatorname{Spec}(A)$ est surjectif.

Ces lemmes étant établis, revenons à la démonstration de l'implication $b)\Rightarrow
a)$ dans (14.4.1). En vertu de (14.3.3.1, (i)), il suffit de prouver que, pour tout
entier $n\geq 0$ et tout point $x'$ de $X'=X[T_1,\ldots,T_n]$ au-dessus de $x$, le
morphisme

